% Fig. 4 -- receiver-conditional observability, and how the four arms isolate the
% paired part of it. Reuses the ellipsoid vocabulary established in Fig. 1.
\begin{tikzpicture}[scale=0.92,every node/.style={font=\footnotesize},
  arw/.style={->,>=stealth,semithick,gray!70}]

% ---- (a) two receivers, two ellipsoids ----
\begin{scope}
  \node[anchor=south,font=\small] at (2.1,2.25) {\textbf{(a)} $\Cobs(\mathcal I_R)$};
  \node (Y) at (0,0.9) {$Y$};
  \draw[arw] (Y) -- (0.95,1.55) node[right,pos=0.95] {$\mathcal I_b=Y$};
  \draw[arw] (Y) -- (0.95,0.25) node[right,pos=0.95] {$\mathcal I_s=(Y,S)$};
  % blind ellipsoid
  \begin{scope}[shift={(3.15,1.62)}]
    \draw[gray!45,dashed] (0,0) circle (0.62);
    \draw[draw=gray!70,fill=gray!12,rotate=20] (0,0) ellipse (0.60 and 0.16);
  \end{scope}
  % side ellipsoid: contains the blind one, never narrower anywhere
  \begin{scope}[shift={(3.15,0.30)}]
    \draw[gray!45,dashed] (0,0) circle (0.62);
    \draw[draw=accent,fill=accent!12,rotate=20] (0,0) ellipse (0.61 and 0.40);
    \draw[draw=gray!70,dotted,rotate=20] (0,0) ellipse (0.60 and 0.16);
  \end{scope}
  \node[anchor=west] at (3.85,1.62) {$\Cobs(Y)$};
  \node[anchor=west,accent] at (3.85,0.30) {$\Cobs(Y,S)$};
  \node[anchor=north,font=\footnotesize] at (2.5,-0.45)
    {$\Cobs(Y,S)-\Cobs(Y)=\Exp[\Cov(m_s\mid Y)]\succeq0$};
\end{scope}

% ---- (b) the four-arm ladder ----
\begin{scope}[shift={(6.9,0)}]
  \node[anchor=south,font=\small] at (1.5,2.25) {\textbf{(b)} isolating the pairing};
  \node (a1) at (0,1.85) {$f(Y)$};
  \node (a2) at (0,1.15) {$f(Y,S_{\mathrm{null}})$};
  \node (a3) at (0,0.45) {$f(Y,S_{\mathrm{shuffled}})$};
  \node (a4) at (0,-0.25) {$f(Y,S_{\mathrm{correct}})$};
  \draw[arw] (a1) -- (a2) node[midway,right=2pt] {architecture};
  \draw[arw] (a2) -- (a3) node[midway,right=2pt] {$S$ marginal};
  \draw[->,>=stealth,very thick,accent] (a3) -- (a4)
    node[midway,right=2pt,accent] {\textbf{pairing}};
\end{scope}
\end{tikzpicture}
